\documentclass{article}
\usepackage{graphicx} % Required for inserting images
\usepackage{hyperref} 

\begin{document}

\section*{Notes}

Taken from introduction

\begin{itemize}
\item In the last ten to fifteen years, the credibility and reliability of quantitative research in both the natural and social sciences has been challenged on several grounds.  Give a brief sketch of the crisis, cite other works that describe it, outline the main dimensions of the crisis.
These problems in quantitative research are mirrored in the public sphere, where constructive discourse is challenged by political polarization and misinformation–a higher order benefit of teaching reproducible quantitative research methods is that students might internalize concepts that make them better/more discerning judges of information they are presented with and more effectively communicate their own arguments/claims/findings (PLAN: move this to the discussion).

\item What can we do as instructors?  We can point out that some progress has been made in the sphere of professional research–like journal data/code policies, pre-registration, results-free refereeing.  There have been some parallel developments in teaching–like Project TIER, Joe and Peter’s work, Nick’s work with Ben Baumer, Mine C-R et al.--but in quantitative methods instruction these issues are still often relegated to the status of special topics and omitted from the core curriculum.  This paper shows some things instructors can do to integrate R&R principles and methods into courses of all types.

\item Fortunately, there are concrete things instructors can do to enhance their classes in ways that address these shortcomings-some are on the mechanical/technical side, some address broader principles/purposes/concepts.  So this paper is meant to help instructors understand those approaches so that they can adopt them in their own classes; our vision is that concepts and principles of reproducibility should become totally standard components of the research methods curriculum across quantitative fields.
\end{itemize}

Notes from 2024-11-25:

- It would be great if some or all of these citations were the kinds of things we would assign to our students. (Accessible, engaging) 
- NASEM, 2019 is a good one
- 1,500 scientists lift the lid on reproducibility

Ioannidis (2001) https://www.nature.com/articles/ng749z

https://journals.sagepub.com/toc/ppsa/7/6


Behavioral Priming: It's All in the Mind, but Whose Mind?
Stéphane Doyen , Olivier Klein, Cora-Lise Pichon, Axel Cleeremans
(2012) https://journals.plos.org/plosone/article?id=10.1371/journal.pone.0029081

Joe's suggestions for background that can be used on syllabi:
- Baker https://doi.org/10.1038/533452a 1,500 scientists lift the lid is nice for its infographics
- The Turing way book: https://book.the-turing-way.org/ 

This one? https://www.nature.com/articles/nature.2015.18248
Maybe this? Schooler, J. Metascience could rescue the ‘replication crisis’. Nature 515, 9 (2014). https://doi.org/10.1038/515009a

https://retractionwatch.com/2016/09/26/yes-power-pose-study-is-flawed-but-shouldnt-be-retracted-says-one-author/

https://retractionwatch.com/the-retraction-watch-leaderboard/top-10-most-highly-cited-retracted-papers/

Notes from 2024-11-14:

\begin{enumerate}
\item we probably want to focus on replicability (and less on reproducibility): but it's important to define, distinguish, put on a continuum
\item introduce replicability and reproducibility in section 1, define in more detail in section 2, define what's intended (replicability vs. reproducibility) in the learning outcomes for each of the sections in 3.1-3.6.
\end{enumerate}

XX example citation \citep{tier}

XX here's another \cite{Sandve:2013} (the latter inline)
XX Joe will draft 3.1 (reproducibility check) 
XX Richard has offered to work on 3.2 research organization

XX Nick will try to iterate on 3.1.

XX Nick and Matt are on the hook for section 3.4
XX Richard will put priority on writing text for the article about the reproducibility classroom activity--worry about completing the exercise after that.

\section{Notes and Outline}

\begin{itemize}
    \item General structure
    \begin{itemize}
        \item Specific examples for a variety of courses
        \item link to github repository
    \end{itemize}
\end{itemize}

\textbf{For Each Vignette}

\begin{enumerate}
\item
    Learning goals and outcomes
\item 
    Scaffolding/skills (emphasize what it paves the way for)
\item 
    Description of teaching 
\item 
    Assessment 
\item 
    Further resources 
\end{enumerate}

\textbf{Five Vignettes}

1. Scott

2. Scott / Joe

3. Nick

4. Matt 

5. Joe / Matt

\begin{enumerate}
\def\labelenumi{\arabic{enumi}.}
\item
  Rotation style final project

  \begin{enumerate}
  \def\labelenumii{\alph{enumii}}
  \item
    Emphasis on final project in SDS 100\\
  \item
    Macalester had a similar course\\
  \item
    This one-credit-hour course uses a module-based approach to
    introduce programming language and principles of reproducibility
    independently from specific course such as introduction to stats,
    etc. This would provide an opportunity for students across an
    institution to learn some fundamental principles without needing to
    commit to a full course. Additionally, this course is competency
    based, meaning students must individually demonstrate competence\\
  \end{enumerate}
\item
  Pre-registration/Plan of Analysis: incorporate in classes.

  \begin{enumerate}
  \def\labelenumii{\alph{enumii}}
  \item
    Useful
    \href{https://www.bitss.org/education/mooc-parent-page/week-2-publication-bias/detecting-and-reducing-publication-bias/p-curve-a-tool-for-detecting-publication-bias/}{video}
    on the implication of p-hacking through p-curve identification\\
  \end{enumerate}
\item
  Teaching Replication Crisis/P-Hacking

  \begin{enumerate}
  \def\labelenumii{\alph{enumii}}
  \item
    Example of using generated data as well as real data to let students
    try for themselves\\
  \end{enumerate}
\item
  Peer Assessment/Self Assessment (sustainable assessment)

  \begin{enumerate}
  \def\labelenumii{\alph{enumii}}
  \item
    Learning goals for peer assessment\\
  \item
    Learning goals for self-assessment\\
  \end{enumerate}
\item
  


  \begin{enumerate}
  \def\labelenumii{\alph{enumii}}
  \item
    Activity where students evaluate papers for reproducibility\\
  \item
    Not necessarily running the replication, but evaluating whether
    there is documentation (code book), data/code available, meta data,
    look at the research design and see if it is specific\\
  \item
    Sources of information: PLOS One, Harvard Data Verse,
  \end{enumerate}
\end{enumerate}


\textbf{Structure of the Paper}

\begin{itemize}
\item
  Introduction

  \begin{itemize}
  \item
    Framing not just for a data analysis course\\
  \end{itemize}
\item
  Principles/Guidelines

  \begin{itemize}
  \item
    Designed for any course, no coding needed, but coding can augment
    these\\
  \item
    Highlight uncertainty in process

    \begin{itemize}
    \item
      There is no right answer\\
    \end{itemize}
  \item
    Active learning\\
  \item
    Scaffolding\\
  \item
    Reflection, ethics, communication\\
  \end{itemize}
\item
  Vignettes

  \begin{itemize}
  \item
    Frameworks

    \begin{itemize}
    \item
      Specific examples for a variety of courses\\
    \item
    \end{itemize}
  \end{itemize}
\item
  Conclusion

  \begin{itemize}
  \item
    Forward-looking\\
  \item
    Returning/synthesis of principles by explicitly tying vignette
    details to principles

    \begin{itemize}
    \item
      Emphasize that this isn't a recipe for specific courses, but
      rather a guide\\
    \end{itemize}
  \item
    Link to Github repository/other place for supplemental course
    materials\\
  \item
    Limitations

    \begin{itemize}
    \item
      Measuring outcomes/success
    \end{itemize}
  \end{itemize}
\end{itemize}

textbf{Other thoughts}

Emily Marshall as reader? (https://sites.google.com/site/emilycorinnemarshall/home/cv?authuser=0) 
Follow up with Javier (Scott) and Richard (Nick)

\end{document}